


血管内介入治疗是实现心血管及脑血管疾病微创诊治的重要手段。在传统手术中，医师需在X射线透视引导下手动推进并旋转导丝与导管，这一过程不仅高度依赖术者的精准操作技巧与临床经验，还使术者长期暴露于电离辐射并承受沉重的铅衣负担\textsuperscript{[1]}。血管介入机器人通过主从遥操作模式驱动柔性器械完成递送与旋转，可将术者与辐射源物理隔离，并提供高精度的可重复运动轨迹。然而，主从遥操作隔绝了术者与器械的直接物理接触，器械在血管内的受力状态必须经过从端测量与主端再现才能传递至术者\textsuperscript{[2]}。因此，在测量器械操作载荷的基础上形成连续、清晰的主端力反馈，是帮助术者判断器械受阻状态的重要环节。

介入器械的力感知方案主要分为体内远端测力和体外近端测力两类，二者反映不同空间尺度的受力信息。远端测力将光纤布拉格光栅（FBG）等微型传感元件集成于器械尖端以直接感知末端接触力\textsuperscript{[3]}，但受限于微细器械尺寸与封装工艺，制造成本高且难以兼容常规一次性耗材。相比之下，近端测力在体外驱动端通过传力路径中的力传感器\textsuperscript{[4,5]}或驱动电机电流\textsuperscript{[6]}获取沿器械传递的综合操作阻力，无需改造常规器械且便于无菌分离，但近端测力通道与机械驱动过程紧密耦合，器械真实负载与传动机构自身的运行干扰将同时作用于测量信号。

器械与驱动机构的接触方式决定了近端传力特性。滚轮摩擦驱动需维持法向预紧力，在推进与旋转复合运动时还会受到界面微滑移的影响\textsuperscript{[6]}；仿人手往复夹持机构则通过面接触建立静摩擦连接\textsuperscript{[7]}，便于控制夹紧力，并分别测量轴向递送力与周向旋转扭矩\textsuperscript{[8]}。受直线导轨行程限制，此类机构通过主、副夹爪“夹持递送---释放回退---重新夹持”的周期性往复运动实现长距离输送\textsuperscript{[9,10]}。对于承担近端测力功能的移动主爪，往复加减速引起的浮动组件惯性、旋转运动引起的支路惯性以及夹爪接触切换会改变测力支路的受力平衡，造成反馈波动。

针对近端测量中的机构附加干扰分量，现有研究主要从运行偏置标定、动力学模型补偿及数据驱动估计等方面展开。在机理建模与标定方面，Sankaran等\textsuperscript{[6]}在不同转速下采集机构自由递送时的空载电流，建立运行偏置与转速的映射关系并从总力中予以扣除。Bao等\textsuperscript{[8]}综合分析了摩擦、重力及旋转速度对测量结果的影响，构建了基于载荷与转速的动态补偿模型，并在时变加载实验中验证了其相较于静态零偏补偿的优势。Wang等\textsuperscript{[11]}建立了递送机构的平移与旋转动力学模型，结合运动学测量实时修正测力基线，并对扭矩信号的抖动进行了滤波处理。在数据驱动建模方面，时序神经网络为表征夹持与传力过程中的复杂非线性关系提供了有效途径。Lyu等\textsuperscript{[12]}将机器人运动状态、夹持间距与软体夹持器的视觉形变序列输入长短期记忆网络（LSTM），实现了器械操作力与扭矩的端到端估计；Wang等\textsuperscript{[13]}采用时间卷积网络（TCN）建立了近端测量与导丝远端接触力之间的时序映射关系。

上述研究分别讨论了运行偏置校正、连续运动补偿和时序力估计，但往复夹持递送还包含主副夹爪交替承载。载荷交接、回程启停惯性以及柔顺接触的变形重建与微滑移，会与真实血管阻力变化同时出现在近端测量信号中，形成卸载台阶和附加波动。因此，需要根据往复夹持的传力过程重构连续递送时的反馈载荷，在减小换手波动的同时保留有效承载阶段新增的血管阻力变化。

\section{问题描述（Problem description）}

\subsection{往复夹持递送过程}


本文研究的夹持式血管介入机器人采用主、副双夹爪交替夹持器械，通过夹爪与器械表面的摩擦接触传递运动和载荷。主夹爪安装于沿直线导轨移动的滑台上，负责器械的轴向推进与周向旋转，其内部集成轴向与周向两路测力支路；副夹爪安装于近血管鞘入口端，在主夹爪释放与回退期间保持器械位置。





受直线导轨行程限制，长距离递送由周期性短行程往复运动完成。根据运动和夹持时序，将一个往复周期划分为S1～S10十个阶段，如图1所示。
\begin{enumerate}
\item 递送与停止（S1～S2）：主爪独立夹持器械推进，并可同时旋转（S1）；单次行程结束时滑台减速停止，副爪尚未闭合（S2）。
\item 卸载交接（S3～S4）：副爪闭合，两爪共同夹持器械（S3）；随后主爪释放，由副爪承载（S4）。
\item 空载回程（S5～S6）：主爪脱离器械后沿导轨退回（S5），在起点减速停止（S6），期间器械由副爪保持。
\item 重新承载与递送（S7～S10）：主爪在起点闭合，两爪再次共同夹持（S7）；副爪释放后，静止的主爪独立承载（S8）；主爪恢复前向推进，轴向载荷逐步恢复（S9）；恢复完成后进入稳定承载阶段（S10）。
\end{enumerate}







JQRBLOCK001






\subsection{测力波动与补偿目标}


往复递送过程中，主爪测力支路经历承载、卸载和重新承载。滑台启停产生轴向惯性力，旋转加减速产生附加力矩，旋转姿态还会改变零载基线。夹爪开合冲击、停止后的载荷松弛及重新推进时的载荷恢复，使近端测量信号出现周期性波动。





本文以相同初始状态、器械位置及推进和旋转轨迹下，主爪持续夹持时的连续递送载荷为参考。在S1末锁存稳态载荷，S2～S8保持该值；S9利用前向推进行程描述轴向载荷恢复，并由GRU网络修正恢复模型的残差。轴向力在有效承载阶段由校正后的实时测量值与补偿量相加得到。扭矩通道进行测量模型校正和S2～S8锁存，在S9恢复输出校正后的测量值。





全文设轴向正方向沿器械名义中心轴指向远端，旋转角正方向遵循右手定则。力、扭矩、位移、角位移和时间的单位分别为 $\mathrm{N}$、$\mathrm{N}\cdot\mathrm{mm}$、$\mathrm{mm}$、$\mathrm{rad}$ 和 $\mathrm{s}$。轴向传感器读数记为 $F_{xm}$，周向支路的切向力读数记为 $F_{tm}$，经标定等效力臂 $\ell_s$ 换算得到测量扭矩 $\tau_{xm}=\ell_sF_{tm}$。输出至主端的反馈信号统一记为 $\bm w_{\mathrm{fb}}=[F_{\mathrm{fb}},\tau_{\mathrm{fb}}]^{\rm T}$。




\subsection{本文工作}


本文围绕往复夹持递送中的轴向力反馈重构开展以下工作。
\begin{enumerate}
\item 建立测量干扰模型和双爪交接受力关系，以S1末稳态载荷和S8末实测载荷确定轴向重加载补偿的初值。
\item 建立以前向推进行程为自变量的轴向载荷恢复模型，结合单输出GRU残差修正、交接锁存和补偿量平滑撤除，形成分阶段反馈方法。
\item 通过参数辨识、配对递送、阻力突变及消融实验评价轴向力重构，并在血管模型中比较操作时间和通过成功率；扭矩通道评价测量校正与锁存效果。
\end{enumerate}





第2节介绍受力关系、轴向载荷恢复模型与反馈方法，第3节给出实验配置和评价方法，第4节与第5节分别为讨论和结论。





\section{往复夹持传力与反馈补偿方法（Load transmission and feedback compensation for reciprocating clamping）}




\subsection{递送测量关系与双爪交接受力}


主爪测力支路与双爪交接关系如图2所示。主、副爪之间的器械段不接触其他物体；副爪远端的管壁接触和实验加载装置对器械的作用统一计入外载。忽略器械自身惯性时，轴向力和绕轴扭矩满足


JQRBLOCK010


式中，$F_M,\tau_M$ 与 $F_S,\tau_S$ 分别为主、副爪作用于器械的轴向推力和驱动扭矩，以远端推进和右手定则正转为正；$F_{\mathrm{out}},\tau_{\mathrm{out}}$ 为外部轴向合阻力和合阻力矩，以阻碍上述正向运动为正，即等于外部作用在统一坐标下分量的相反数。记 $\bm w_M=[F_M,\tau_M]^{\rm T}$，连续递送参考载荷为 $\bm w_{\mathrm{ref}}=[F_{\mathrm{ref}},\tau_{\mathrm{ref}}]^{\rm T}$。副爪完全释放时，$F_S=\tau_S=0$，主爪平衡当时的全部外载。





将轴向力与换算后的扭矩组成测量向量 $\bm y=[F_{xm},\ell_sF_{tm}]^{\rm T}$。传感器连接浮动传力组件与滑台；内侧浮动组件的轴向弹性位移较小，其绝对加速度近似为滑台加速度。测量关系写为


JQRBLOCK011



JQRBLOCK012


式中，$\bm q_M=[x_M,\dot x_M,a_M,\varphi_M,\dot\varphi_M,\alpha_M]^{\rm T}$ 为平移轴和旋转轴的位置、速度及加速度；$m_s$ 和 $J_s$ 为两支路的等效质量与转动惯量，其单位与式(3)中的毫米制运动量相匹配。$\bm H_0\in\mathbb R^{2\times2}$ 为离线标定的可逆传递矩阵，交叉通道元素包含力与力矩之间的单位换算。$\bm b(\varphi_M)$ 为包含静态零偏和角度相关重力分量的零载基线，$\bm\varepsilon$ 为测量噪声。各向量的两个载荷分量依次以 $\mathrm{N}$ 和 $\mathrm{N}\cdot\mathrm{mm}$ 计。





开合指令记为 $\bm c=[c_M,c_S]^{\rm T}\in\{0,1\}^2$，闭合取1、释放取0。$\bm p_{\mathrm{clamp}}(t)$ 表示开合动作触发的短时冲击，随动作后的时间变化。开合动作位于S3、S4、S7和S8，反馈在这些阶段采用锁存值，冲击项用于说明测量组成，不参与在线解析补偿。





S3和S7中两爪同时闭合，夹持垫与传力支路均具有柔顺性。新闭合的夹爪在尚未建立明显切向变形时，载荷仍主要由原承载爪承担。取S7末和副爪完全释放后的S8末，由式(1)得到


JQRBLOCK013


上标 $(7)$、$(8)$ 表示对应阶段末端。主爪载荷增量取决于副爪原先承担的力及交接前后的外载变化；若外载不变，主爪接管副爪原先承担的全部轴向力。S8末主爪读数低于换手前稳态值，表示当前实际载荷与连续递送参考之间存在差值。





S8末副爪已完全释放，开合冲击已衰减，平移轴和旋转轴均静止，此时 $a_M=\alpha_M=0$。取该有效采样时刻为 $k_9$，载荷初值为


JQRBLOCK014









JQRBLOCK002







\subsection{换手后重新推进阶段的轴向载荷恢复模型}





先定义扣除惯性与零载基线后的测量值


JQRBLOCK015


在S1匀速段末锁存连续递送的局部参考值


JQRBLOCK016


式中，$k_{\mathrm{S1,end}}$ 为减速指令发出前的最后一个采样时刻，$N_{\mathrm{lock}}$ 为完全位于匀速段内的窗口长度。轴向初始载荷缺口为


JQRBLOCK017







轴向恢复模型用于前向推进且 $\widehat A_{9,F}>0$ 的过程。从S8末起累计前向推进行程


JQRBLOCK018


本文将恢复初期的载荷增长视为传力链的弹性恢复，并假设停顿后的附加摩擦阻力随推进行程衰减。令 $w_F(s)$ 表示模型给出的轴向载荷，取 $s=s_{\mathrm{rec}}$，采用一阶恢复关系


JQRBLOCK019


式中，$\lambda_F>0$ 为弹性恢复的特征行程，$\Delta_F\ge0$ 为停顿后附加阻力幅值，$D_c>0$ 为该附加项衰减的特征行程。以 $w_F(0)=\widehat w_{M,k_9,F}$ 为初值，积分得到相对于锁存参考值的补偿量


JQRBLOCK020


当 $D_c=\lambda_F$ 时，采用极限形式


JQRBLOCK021


当 $\Delta_F=0$ 时，式(11)退化为单指数恢复。双指数形式同时描述初始载荷增长和附加阻力衰减；出现高于锁存参考值的峰值时，$f_{\mathrm{rel},F}$ 为负，将该附加分量从实测值中扣除。





恢复起点的加载斜率定义为


JQRBLOCK022


在弹性恢复假设下，待恢复载荷越大、初始加载斜率越小，对应的恢复特征行程越长。$k_0$ 由恢复曲线确定，跨周期比较 $\lambda_F$ 与 $\widehat A_{9,F}+\Delta_F$ 的比例关系，可检验初始加载斜率是否近似稳定。





令式(10)右端为零，得到有限峰值对应的行程


JQRBLOCK023


对于 $\widehat A_{9,F}>0$、$\Delta_F>0$，$D_c>\lambda_F$ 时存在有限峰值；$D_c<\lambda_F$ 时还需满足 $\widehat A_{9,F}<\Delta_FD_c/(\lambda_F-D_c)$。两尺度相等时，$s^*=\lambda_F(1+\widehat A_{9,F}/\Delta_F)$；$\Delta_F=0$ 时单调趋近参考载荷。





参数由离线重加载曲线辨识。有峰曲线存在等效的尺度交换形式，统一以 $0<\lambda_F\le D_c$ 的分支报告，跨周期比较使用同一分支；单调曲线按其实际拟合结果报告。用于反馈验证的参数仅由辨识集确定，验证前固定；逐周期拟合用于分析恢复形状及初始斜率。恢复过程及补偿项如图3所示。







JQRBLOCK003







\subsection{GRU轴向残差修正与分阶段反馈}





轴向物理基线由校正实测值与恢复补偿量组成：


JQRBLOCK024


其中，$\mathbb I_{\mathrm{rec}}=1$ 表示S9中已启用且尚未开始撤除的轴向恢复补偿，其余情况为0。补偿启用时，物理基线在恢复起点满足 $F_{\mathrm{base}}[k_9]=\widehat D_{9,F}$。





GRU网络用于估计恢复模型未描述的轴向残差。输入向量为


JQRBLOCK025


式中，$\Delta t[k]$ 为距S9起点的时间，$L_{\mathrm{ins}}[k]$ 为扣除空载回程后的累计有效插入深度，用于描述随插入位置变化的恢复残差。输入中的载荷量在本次S9开始时确定，网络不接收S9的实时测力值。





输入窗口为 $\bm X_k=[\bm x_{k-N+1},\ldots,\bm x_k]^{\rm T}$。每次S9开始时重置窗口，以首个有效样本补齐不足的历史长度。网络由输入标准化层、两层单向GRU及单输出线性层组成，采用标准GRU更新关系\textsuperscript{[14]}，反标准化后得到轴向残差 $e_F[k]$。残差的逐点限幅及淡入为


JQRBLOCK026


式中，$\operatorname{sat}(e,a)=\max(-a,\min(e,a))$，$\eta\ge0$ 为限幅比例，$e_{0,F}>0$ 为最小限幅，$N_b$ 为淡入采样点数。仅在 $\widehat A_{9,F}>0$ 时启用恢复补偿和残差修正。





在正常恢复阶段，总补偿量为 $C[k]=f_{\mathrm{rel},F}(s_{\mathrm{rec}}[k])+r_F[k]$，反馈为


JQRBLOCK027


补偿启用时，S9起点 $\beta_{k_9}=0$，故 $F_{\mathrm{fb}}[k_9]=\widehat D_{9,F}$，与交接锁存值衔接。正常恢复期间，实测值以单位权重进入输出。





恢复完成行程定义为


JQRBLOCK028


该判据取补偿量最终进入误差带的行程。当 $s_{\mathrm{rec}}\ge s_{\mathrm{th}}$，或前向速度低于 $v_{\min}$ 持续超过 $T_{\mathrm{st}}$，或出现 $\dot x_M<-v_{\min}$ 的回拉时，开始撤除补偿。记触发时刻为 $t_e$，保存触发前最后有效补偿量 $C_e$，随后输出


JQRBLOCK029


式中，$v_{\min}$ 为最小有效前向速度，$T_{\mathrm{st}}$ 为停止判定时长，$T_{\mathrm{off}}$ 为补偿撤除时长。撤除期间停止更新恢复模型与GRU输出，$\rho_k=0$ 后仅输出校正实测值；因停止或回拉而退出后，本次恢复补偿不再重新启用。正常恢复完成并撤除补偿后进入S10。





令 $C_{\mathrm{out}}[k]$ 在正常恢复时为 $C[k]$，撤除时为 $\rho_kC_e$，未启用或已撤除时为0。全周期反馈写为


JQRBLOCK030


当 $\widehat A_{9,F}\le0$ 时，S9的 $C_{\mathrm{out}}=0$。扭矩在S9直接恢复为校正实测值。





网络训练采用与在线淡入和限幅一致的输出。令 $\Omega_{\mathrm{S9}}$ 为训练集中启用正常恢复补偿的S9样本，$\mathcal P_{\mathrm{S9}}$ 为同一恢复段内、均属于 $\Omega_{\mathrm{S9}}$ 且时间相邻的样本对，损失为


JQRBLOCK031


式中，$F_*>0$ 为由训练集确定的力归一化尺度，$\gamma_{\mathrm{sm}}$ 为平滑权重。训练、验证与测试按完整试次划分，输入标准化参数仅由训练集计算。方法流程如图4所示。







JQRBLOCK004






\section{实验验证与分析（Experimental validation and analysis）}

\subsection{实验平台}


实验平台基于双操作手协同血管介入机器人构建，受测器械包括0.014英寸导引导丝与5F介入导管，主夹爪滑台平移及旋转由伺服电机驱动。近端测力机构采用两只力恒LH-S09A3单轴拉压传感器，分别测量轴向力和切向力；切向力经标定等效力臂换算为扭矩。测量与控制系统基于倍福控制总线构建，以 $250\,{\rm Hz}$ 同步记录测力读数、各轴位置、速度、加速度及夹爪开关指令。数据处理在工控机（Intel Core i7-11700，16 GB RAM）上完成。





实验配置分为以下两类，如图5所示。
\begin{enumerate}
\item 测量参数标定配置。将测力支路与外部标准力传感器串联，施加已知轴向载荷及已知力臂的切向载荷，标定传递矩阵 $\bm H_0$ 和传感器等效力臂 $\ell_s$。空载角度扫描与平移、旋转启停试验分别用于辨识零载基线和惯性参数。
\item 配对递送与血管模型实验配置。将器械置入血管模型或具有远端加载单元的管路，分别执行换手递送和不换手连续递送。换手试次从给定起点推进 $8\sim10\,{\rm mm}$ 后完成换手，再推进约 $10\,{\rm mm}$；配对试次从同一起点连续推进相同总行程。两类试次采用相同的器械、初始受力状态、管路和递送轨迹，推进阶段的速度及旋转轨迹保持一致。
\end{enumerate}





配对试次按累计有效推进位移对齐，空载回程不计入有效推进位移；伴随旋转的试次同时匹配旋转轨迹。以不换手试次经惯性和零载基线校正的载荷作为 $\bm w_{\rm ref}$。在交接停止区间，参考值取对应有效推进位置的连续递送载荷。重复不换手试次用于估计参考载荷的重复性。具体管路、远端加载方式、试次分组及重复次数为[待补]。







JQRBLOCK005







\subsection{测量参数标定与载荷恢复模型辨识}





通过静态加载得到 $\bm H_0$ 与 $\ell_s$，通过空载角度扫描得到 $\bm b(\varphi_M)$，通过轴向与旋转启停数据分别辨识 $m_s$ 和 $J_s$。各参数、标定范围和残差统计为[待补]。





为观察夹爪开合时的承载变化，在交接试次中适当延长S3与S7的保持时间。该组记录中，静止状态下夹爪闭合产生约 $15\sim20\,{\rm ms}$ 的短暂波动，随后无明显轴向载荷转移；S4主爪释放和S8副爪释放分别对应主爪卸载与重新承载。S8末主爪载荷增量占S7末副爪载荷的比例约为 $0.74$。该比例及其测量来源、重复试次统计为[待补]。





对S9轴向恢复曲线辨识 $\lambda_F$、$\Delta_F$ 和 $D_c$，并由式(13)计算初始加载斜率 $k_0$。一组恢复曲线的初始加载斜率约为 $1.28\,{\rm N/mm}$。所观察的典型曲线在约 $3.7\,{\rm mm}$ 处出现峰值，峰后载荷回落约 $0.22\,{\rm N}$；相应曲线、样本范围和统计结果为[待补]。峰值位置及回落幅值随推进速度、停留时间和试验周期的变化另作统计，$D_c$ 的跨周期变化结果为[待补]。





采用相同的辨识区间分别拟合单指数模型（$\Delta_F=0$）与双指数模型，比较其对载荷上升、峰值和峰后回落的描述。除整个S9区间的拟合RMSE外，分别报告前 $1\,{\rm mm}$ 和前 $2\,{\rm mm}$ 的拟合RMSE，以量化部分周期起步段的偏差，结果为[待补]。拟合误差以校正后的实测恢复曲线为基准；连续力反馈误差按第3.3节的配对参考计算。





跨周期关系检验包括两部分。首先，将 $\lambda_F$ 对 $\widehat A_{9,F}+\Delta_F$ 作回归，结合各周期 $k_0$ 的分布，检验两者是否近似成正比。其次，对存在有限峰值的周期，按式(14)计算峰值行程并与实测值比较；单调恢复的周期单独报告其数量。回归系数、拟合区间、峰位误差和对应图表为[待补]。





物理参数辨识、网络训练和独立测试按完整试次划分，测试前冻结模型参数及标准化参数。逐周期拟合用于分析参数与曲线特征的关系。GRU采用2层单向结构，隐层规模分别为64和32，输入窗口长度为30个采样点，对应约 $0.12\,{\rm s}$；输出为一个轴向残差。训练、验证与测试试次数，优化参数，锁存窗口 $N_{\rm lock}$、残差淡入长度 $N_b$、限幅参数 $\eta$ 与 $e_{0,F}$、停止判定参数及撤除时长 $T_{\rm off}$ 为[待补]。




\subsection{对比实验结果与指标分析}


各方法采用同一试次预先确定的阶段边界和评价窗口。轴向力比较采用以下六种方法。Raw仅作传递矩阵换算，保留原始波动；Kinematic Bias在此基础上扣除惯性和零载基线；ZOH在S2至S9保持S1末锁存值，进入S10后输出校正实测值；Physics Baseline采用轴向恢复模型及本文的阶段切换规则，不加入GRU；Physics + Coupled GRU在相同补偿流程中，将实时轴向物理基线额外输入GRU；Physics + Decoupled GRU采用本文不含实时测力输入的残差网络。后两种网络采用相同的数据划分、网络规模、限幅和输出过渡设置。扭矩单独评价测量校正与S2～S8锁存的效果，S9直接使用校正实测值。





\subsubsection{配对参考误差分析}





对稳态递送阶段（S1、S10）和轴向重加载阶段（S9）分别计算均方根误差（RMSE）与平均绝对误差（MAE）。令 $d_p[k]=w_{{\rm fb},p}[k]-w_{{\rm ref},p}[k]$，则


JQRBLOCK032


式中，$\Omega_s$ 为相应阶段的样本集合，$p\in\{F,\tau\}$ 为评价通道。交接锁存阶段另行报告误差，不并入稳态递送阶段。在S2～S10的评价窗口内，误差峰峰值定义为


JQRBLOCK033







轴向反馈恢复时间采用绝对误差带定义：


JQRBLOCK034


式中，$t_{\rm win}$ 为连续保持在误差带内的时间，$\varepsilon_{\rm tol}=2\sigma_{{\rm rep},F}$。$\sigma_{{\rm rep},p}$ 由重复不换手试次估计：在各对齐位置计算试次间的载荷方差，再取评价区间内方差均值的平方根。两通道的重复性分别以力和扭矩单位报告。观察区间内未满足恢复条件的试次记为未恢复，并列出其数量。误差保持时间与观察区间长度为[待补]。







JQRBLOCK009







表1按器械和递送工况汇总轴向结果，各方法的重复试次统计及差异分析为[待补]。扭矩通道报告校正前后的误差、交接锁存期间的输出变化及S9恢复实测输出时的变化，结果与参考重复性 $\sigma_{{\rm rep},\tau}$ 为[待补]。




\subsubsection{阻力突变响应评估}


在S1稳态递送中途和S9初始恢复阶段，通过远端加载单元施加未参与训练的轴向阶跃阻力。换手试次及配对不换手试次在相同递送位置施加相同扰动，并设置相应的无扰动重复试次。用有、无扰动波形之差提取阻力增量，分别记为 $\delta F_{\rm fb}(t)$ 和 $\delta F_{\rm ref}(t)$。阶跃保留比与响应延迟定义为


JQRBLOCK035



JQRBLOCK036


其中，$\Delta F$ 为增量波形在规定阶跃平台内的均值，$t_{{\rm fb},63\%}$ 与 $t_{{\rm ref},63\%}$ 均从各自试次的阶跃施加时刻开始计时，分别为增量波形首次达到参考阶跃幅值 $\Delta F_{\rm ref}$ 的63\%所需的时间。未达到该阈值的试次记为未响应。阶跃幅值、施加位置、平台窗口及统计结果为[待补]。重点比较耦合输入与解耦输入GRU在S9中的幅值保留和响应延迟。







JQRBLOCK006






\subsubsection{消融实验与模型泛化性分析}


以完整方法为基准，分别设置去掉轴向恢复模型、去掉附加衰减项（$\Delta_F=0$）和去掉GRU的消融变体。各变体采用相同数据划分和评价窗口，含GRU的变体按各自物理基线重新训练。比较S9误差、前 $1\sim2\,{\rm mm}$ 的误差及全周期波动，结果见图7的预留位置。





泛化性评价包括不同推进速度、夹爪预紧力及器械条件。推进速度设置为 $5\,{\rm mm/s}$、$10\,{\rm mm/s}$ 和 $15\,{\rm mm/s}$，预紧力范围和试次分配为[待补]。各测试条件使用冻结参数，分别报告轴向恢复误差和阶跃响应。另记录S9完成、停止和回拉时的补偿撤除波形，以及初始载荷缺口非正时的直接测量输出；参数设置与结果为[待补]。







JQRBLOCK007







\subsubsection{血管模型连续递送重构与操作者感知实验}





在介入血管模型中执行多周期往复推进及伴随旋转的递送任务，记录原始测量、校正测量和反馈输出。轴向波形用于评价换手后的连续力反馈重构，扭矩波形用于观察测量校正、交接锁存及S9直接输出的变化，图8预留完整记录及局部放大结果。





操作者实验采用主从递送任务，对比未补偿、传统测量校正和本文轴向连续力反馈三种条件。各条件使用相同的扭矩处理和主端反馈映射，比较因素为轴向反馈方式。由多名操作者沿指定路径递送器械，并完成狭窄或大曲率分叉等高难部位的通过任务。操作者人数、经验分组、熟悉训练、条件执行顺序及每条件重复次数为[待补]。





介入总时间从任务开始计至器械到达目标位置，通过成功率为成功完成指定高难部位的试次数占该部位总试次数的比例。目标位置、成功判定、失败条件和最长允许时间为[待补]；未完成任务单独报告。按操作者配对汇总各条件的完成时间与成功率，统计方法、结果及差异分析为[待补]。







JQRBLOCK008






\section{讨论（Discussion）}


轴向恢复模型将初期载荷增长与随后附加阻力的衰减合并描述。初始加载斜率 $k_0$ 联系了载荷缺口、附加阻力幅值与恢复特征行程；在弹性恢复假设下，较低的加载斜率对应较长的恢复行程。跨周期比例关系及峰位比较结果为[待补]。部分周期起步段与物理模型的偏差由GRU残差修正，其贡献通过初始行程误差和消融结果评价。





实测值与补偿量分别进入轴向输出，使恢复模型和GRU的过渡处理集中于补偿量。阶跃实验评价新增阻力的幅值保留与响应延迟，血管模型实验进一步比较不同轴向反馈条件下的介入总时间及高难部位成功率；相关结果及其对操作者判断受阻状态的影响为[待补]。





交接锁存期间（S2～S8）的适用条件为：减速停止后的器械由当前承载夹爪保持，连续递送参考载荷在该短时段内近似不变。扭矩在S9采用校正实测值，长柔器械在弯曲管腔中的扭转储能与甩转未纳入轴向恢复模型。




\section{结论（Conclusions）}


本文针对往复夹持递送引起的近端测力波动，建立了测量校正、交接受力分析、轴向重加载模型与GRU残差修正相结合的反馈方法。以S1末稳态载荷和S8末实测承载确定轴向补偿初值，用前向行程描述弹性恢复及附加阻力衰减，并在恢复完成、停止或回拉时平滑撤除补偿。扭矩通道采用测量校正与交接锁存。配对递送误差、阶跃响应、消融及操作者实验的定量结论为[待补]。




\noindent
\textbf{数据可用性声明}


[待补：与本研究对应的数据集名称、存储平台、访问链接及开放状态。]








JQRBLOCK037



\section*{作者简介}
 第一作者：隋凯丞（2000--），男，硕士研究生。研究领域：血管介入手术机器人，机器人力感知与控制。E-mail: suikaicheng24@mails.ucas.ac.cn。



通信作者：胡琨元（出生年份待补），作者简介与E-mail为[待补]。




 通信作者：刘浩（1976--），男，博士，研究员，博士生导师。研究领域：医疗微创手术机器人，智能控制。E-mail: liuhao@sia.cn。



